\begin{problema}\leavevmode
\begin{enumerate}[label=\textbf{\alph*)}]
    \item Usando la regla del trapecio compuesta, para intervalo uniformemente espaciados, tenemos que
    \begin{align*}
    I = \int_{0}^{1} f(x) \, dx &\approx 0.25 \left[ \frac{1}{2}\left(f(0)+f(1)\right) + f(0.25) + f(0.5) + f(0.75) \right] \\[1em]
    &= 0.25 \left[ \frac{1}{2}\left(1+0.6703\right) + f(0.25) + 1.0148 + 1.0063 + 0.3819 \right] \\[1em]
    &= 0.934538
    \end{align*}
    
    Similarmente, usando la regla de Simpson 1/3 compuesta para \(n=4\) tenemos
    \begin{align*}
        I = \int_{0}^{1} f(x) \, dx &\approx \frac{0.25}{3}\bigg[f(0)+4f(0.25)+2f(0.5)+4f(0.75)+f(1)\bigg]\\[1em]
        &=\frac{0.25}{3}\bigg[1+(4\times 1.0148)+(2\times 1.0063)+(4\times 0.8819)+0.6703\bigg]\\[1em]
        &=0.939142
    \end{align*}

    \item Estimaremos los errores con ayuda de wolfram mathematica. Sea \(h=0.25\) y \(x_n=nh\) para \(0\le n\le 4\). Utilizaremos diferencias finitas centradas en los puntos interiores \(x_1,\ x_2\) y \(x_3\) y diferencias finitas hacia adelante y hacia atrás para los puntos $x_0$ y $x_4$ respectivamente. Una vez aproximados los valores para $f'(x_n)$, aproximaremos $f''(x)$ siguiendo la misma estrategia, y asi sucesivamente para las demás derivadas de orden superior.
    
    Tenemos entonces que
\
\begin{table}[h]
    \centering
\begin{tblr}{L|L|L|L|L|L}
x                            & 0 & 0.25 & 0.5 & 0.75 & 1 \\[1em] \hline 
f(x)                         & 1 & 1.0148 & -1.0063 & = 0.8819 & 0.6703 \\[1em] \hline 
f'(x)                        & 0.0592 & 0.0126 & -0.2658 & -0.672  & -0.8464 \\[1em] \hline 
f''(x)                      & -0.1864 & -0.65 & \SetCell[c=1]{c, orange} -1.3692 & -1.1612 & -0.6976 \\[1em] \hline 
f'''(x)                      & -1.8544 & -2.3656 & -1.0224 & 1.3432 & 1.8544 \\[1em] \hline 
f^{(iv)}(x)                  & -2.0448 & 1.664 & \SetCell[c=1]{c, orange} 7.4176 & 5.7536 & 2.0448 
\end{tblr}
\end{table}

Hemos resaltado en naranja la segunda y cuarta derivada de valor absoluto más grande. Notemos que en este caso $b=1,\ a=0$ y $n=4$. Entonces 

\begin{align*}
    |E_T|&\approx\frac{1.3692}{12\times 4^2} & |E_S| &\approx \frac{7.4176}{180\times 4^4}\\[1em]
     &\approx 0.00713125 & &\approx 0.000160972
\end{align*}

\item Dado el valor  de referencia $I_{ref}=0.92782$ primero redondeamos el valor aproximado de $I$ obtenido con cada regla, al mismo número de cifras significativas que $I_{ref}$. Se tiene entonces
\begin{align*}
    I_T&=0.93454 \qquad \text{Regla del trapecio}\\[1em]
    I_S&=0.93914\qquad \text{Regla del Simpson 1/3}
\end{align*}

Los errores relativos porcentuales serán, respectivamente para la regla del trapecio y de Simpson
\begin{align*}
    E_T=0.72428\% && E_S=1.2206\%
\end{align*}

Vemos entonces, que para los datos dados, el método del trapecio resulta ser más eficaz que la regla de Simpson en este caso, al tener un error relativo menor. Esto contras con las aproximaciones para el error obtenidas en el inciso anterior, las cuales sugieren lo contrario. Una posible explicación para esto es que, al calcular hasta la cuarta derivada mediante diferencias finitas, en los puntos dados, los errores se acumulan, por lo que las cotas, en particular para la regla de Simpson, resultan no ser muy exactasS.

\end{enumerate}
\end{problema}